% Supplementary Data for the NARGAB manuscript. Compiles on its own.
\documentclass[11pt]{article}
\usepackage[a4paper,margin=2.4cm]{geometry}
\usepackage[T1]{fontenc}
\IfFileExists{lmodern.sty}{\usepackage{lmodern}}{}
\usepackage{amsmath,amssymb}
\usepackage{booktabs,tabularx}
\usepackage{graphicx,xcolor}
\usepackage[numbers,sort&compress]{natbib}
\usepackage[hidelinks]{hyperref}
\usepackage{caption}
\captionsetup{font=small,labelfont=bf}
\newcommand{\todo}[1]{\textcolor{red}{[TODO: #1]}}
\newcommand{\figbox}[4]{%
  \fbox{\parbox[c][#2][c]{0.94\linewidth}{\centering\small
  \textbf{#1}\\[4pt]\begin{minipage}{0.86\linewidth}\raggedright #3\\[3pt]
  \textit{Data:} \texttt{#4}\end{minipage}}}}
\newcommand{\Ecoli}{\textit{E.~coli}}
\newcommand{\nap}{nAP}
\renewcommand{\thesection}{S\arabic{section}}
\renewcommand{\thetable}{S\arabic{table}}
\renewcommand{\thefigure}{S\arabic{figure}}
\renewcommand{\theequation}{S\arabic{equation}}

\title{Supplementary Data\\[2pt]
\large Genomic antimicrobial resistance prediction under lineage and cohort shift}
\author{}
\date{}

\begin{document}
\maketitle

\section{Supplementary Methods: the TriLens architecture}

This section gives the full specification summarised in the main text. Each component produces a
raw score $z$, converted to a probability by the logistic sigmoid
\begin{equation}
\sigma(z) = \frac{1}{1 + \exp(-z)}.
\label{eq:sigmoid}
\end{equation}

\subsection{Arrays passed to the model}
Preprocessing produces six arrays over the 6,962 strains: SNP blocks by gene
($6{,}962 \times 5{,}000 \times 64$, mutation patterns inside genes); gene presence/absence
($6{,}962 \times 32{,}441$, genes acquired from elsewhere); per-gene statistics
($6{,}962 \times 5{,}000 \times 4$, coarse gene-level variation); a phylogroup one-hot
($6{,}962 \times 8$); the label matrix ($6{,}962 \times 7$); and its mask
($6{,}962 \times 7$). The collection is wide and short --- more features than strains even after
reduction --- and patchy, with coverage from 49\% to 87\% and prevalence from 0.12 to 0.50 by drug,
which is why no averaged score over the panel is reported without the per-drug values beside it.

\subsection{Gated attention}
All three branches face the same problem: a strain carries thousands of genes, few matter for any
one drug, and which ones matter changes with the drug. Each therefore pools with the gated attention
of Ilse et al.\ \cite{ilse2018attention}, an attention variant \cite{vaswani2017attention} for bags
in which few items are informative. Each element $h$ receives a relevance score
\begin{equation}
e = w^\top \big(\tanh(V h) \odot \sigma(U h)\big),
\label{eq:gated_attn}
\end{equation}
with $U$, $V$ and $w$ learned and $\odot$ the elementwise product. The two branches serve different
purposes: $\tanh$ proposes how valuable the element looks, in $[-1, 1]$, and $\sigma$ gates how much
of that proposal passes, in $[0, 1]$, so an element that looks informative can still be suppressed
--- which ordinary attention cannot do. Scores are normalised by a softmax and the pooled summary is
the weighted sum of elements.

\subsection{Branch A: the motif encoder}
Branch A reads the SNP blocks, of shape $(G, L)$ with $G = 5{,}000$ genes and $L = 64$ positions.
It looks for short local motifs of mutations inside a gene, on the reasoning that one mutation is
usually silent while two or three close together are more likely to alter the protein. Two stacked
convolutions run inside each gene,
\begin{equation}
h_g = \mathrm{ReLU}\!\left(\mathrm{Conv1d}_{16 \to 32}\big(\mathrm{ReLU}(\mathrm{Conv1d}_{1 \to 16}(x_g))\big)\right),
\label{eq:branchA}
\end{equation}
with windows of width three and dropout between layers, so the first layer learns 16 simple patterns
and the second 32 compound ones. Masked max-pooling then reduces each gene to a fixed-length vector
by keeping, per learned pattern, the strongest response anywhere in that gene, with padding
positions excluded. The per-gene vectors are pooled across genes by \eqref{eq:gated_attn} to give
$z_A \in \mathbb{R}^{64}$ and per-drug scores $\ell_A \in \mathbb{R}^{7}$.

Confining convolution to gene blocks follows a measured result rather than an argument. A plain
one-dimensional convolutional network over the whole pan-genome variant sequence,
\begin{equation}
h^{(l)} = \mathrm{MaxPool}\!\left(\mathrm{ReLU}\!\left(\mathrm{BN}\!\left(\mathrm{Conv1d}_l\!\left(h^{(l-1)}\right)\right)\right)\right),
\label{eq:naive_cnn}
\end{equation}
with four layers of widths 9, 7, 5, 3 and channel sizes 64, 96, 128, 128, the last reduced by global
average pooling and a linear readout, scored 0.775 average precision for ciprofloxacin on the
training collection and 0.632 after the variant order was shuffled, a drop of 0.143; on the Giessen
collection, whose variants are ordered along a single reference, the drop was 0.180. Pan-genome
column order is therefore not meaningless, but it is also not the biological locality a convolution
assumes, and within a gene sequence order is real. Branch A convolves only there.

\subsection{Branch B: the gene attention module}
Branch B reads the presence/absence profile, a binary vector over 32,441 loci, to catch resistance
conferred by an arriving gene rather than a mutation. Each locus has a learned embedding, allowed to
act only if the strain carries it,
\begin{equation}
\tilde{e}_k = x_k \, E_k, \qquad k = 1, \dots, G_{\mathrm{pa}},
\label{eq:branchB_embed}
\end{equation}
where $E_k \in \mathbb{R}^{d_{\mathrm{token}}}$ is the embedding of locus $k$ and
$x_k \in \{0, 1\}$ its presence. Absent loci contribute the zero vector, and their attention scores
are additionally set to $-\infty$ before the softmax, which drives their weight to exactly zero and
reaches the same conclusion twice.

\paragraph{Why \eqref{eq:gated_attn} cannot be applied directly.} Gated attention scores each item
from its own representation. For any locus the strain carries, that representation is
$\tilde{e}_k = E_k$ --- the same vector in every strain --- so the score is strain-independent, the
softmax merely renormalises over whichever loci the strain happens to carry, and the branch
collapses to a global ranking of loci by presence. It cannot express that a locus matters for a
particular isolate, which is the whole point of attending here. The other two branches do not have
this defect, because their inputs already vary by strain. This is the only place in the design where
we alter the mechanism rather than the data it reads.

\paragraph{Strain conditioning.} We condition the score on the mean embedding of the loci the strain
actually carries,
\begin{equation}
c = \frac{1}{\max(n_b, 1)} \sum_{k=1}^{G_{\mathrm{pa}}} \tilde{e}_k,
\qquad n_b = \sum_k x_k,
\label{eq:branchB_ctx}
\end{equation}
so that each locus is scored in the company it keeps,
\begin{equation}
s_k = w^\top\Big(\tanh\big(V \tilde{e}_k + V_c\, c\big) \odot
\sigma\big(U \tilde{e}_k + U_c\, c\big)\Big),
\label{eq:branchB_score}
\end{equation}
with $V_c$ and $U_c$ added projections. A sum of two projections equals one projection on the
concatenation $[\tilde{e}_k;\, c]$, but broadcasting $c$ avoids realising a tensor of size
$B \times 32{,}441 \times 2 d_{\mathrm{token}}$, which at this number of loci is the difference
between affordable and not; the extra cost is two
$d_{\mathrm{token}} \times d_{\mathrm{attn}}$ matrices. The attention sum over 32,441 loci is
accumulated in single precision, because at half precision the round-off is large enough to disrupt
training. The attended embeddings are pooled and projected,
\begin{equation}
z_B^{\mathrm{pre}} = \sum_{k=1}^{G_{\mathrm{pa}}} a_k\, \tilde{e}_k,
\qquad
z_B = \mathrm{GELU}\!\left(W_B\, z_B^{\mathrm{pre}} + b_B\right) \in \mathbb{R}^{64},
\label{eq:branchB_out}
\end{equation}
with $\ell_B \in \mathbb{R}^{7}$ read off $z_B$ by seven independent linear heads. The main text
ablates \eqref{eq:branchB_ctx}--\eqref{eq:branchB_score} and reports the result as a trade rather
than an improvement.

A dense layer over 32,441 inputs would need tens of thousands of weights per unit against about
5,000 training strains and would overfit immediately; attention keeps the parameter count fixed as
the number of loci grows, since the same small scoring function is applied to each locus, and its
weights are readable, which is what the interpretability analysis in the main text reads.

\subsection{Branch C: the cross-antibiotic module}
Branch C reads the per-gene statistics, of shape $(G, 4)$, and exploits relations between the drugs
rather than extracting a new kind of genomic signal. Two drugs of one class are often resisted by
one mechanism, so evidence from a well-represented drug can help a poorly represented one; with
seven independent models, gentamicin, which has few resistant strains, has little to learn from.

Our own data support that reasoning only in part (Figure~\ref{fig:coresistance}). The two
cephalosporins are by far the most strongly associated pair, at $\phi = 0.84$, and nothing else
approaches it --- which is the shared-mechanism account exactly. Below that pair, class membership
stops predicting association: ampicillin pairs with trimethoprim--sulfamethoxazole at 0.55, across
two unrelated classes, as strongly as with amoxicillin--clavulanic acid at 0.54 within its own, and
ciprofloxacin, which has no same-class partner, sits at 0.51 and 0.52 with the two cephalosporins.
Resistance genes on one plasmid or integron travel together whether or not they target the same
drug, and this collection shows both patterns at once. The values are also not a ranking: some pairs
are measured on 235 strains (AMC and SXT) and others on 6,030 (CIP and GEN).

That shapes the design. A module exchanging information only within classes would catch the
cephalosporin link and no cross-class link, so the class-specific heads are followed by a refinement
step free to learn relations the hierarchy does not contain. Each gene's four statistics pass
through a small per-gene perceptron, gated attention pools across genes to
$z_C \in \mathbb{R}^{64}$, and the pooled representation splits into one projection per drug class
(penicillins AMP and AMC; cephalosporins CTX and CAZ; fluoroquinolone CIP; aminoglycoside GEN;
folate-pathway inhibitor SXT),
\begin{equation}
t = \mathrm{GELU}(W_t z_C + b_t), \qquad
u_g = \mathrm{GELU}(W_g t + b_g),
\label{eq:branchC_trunk}
\end{equation}
with $t \in \mathbb{R}^{32}$ shared across all seven drugs and $u_g \in \mathbb{R}^{16}$ specific to
class $g$. Each drug reads an initial score from its own class projection,
\begin{equation}
\ell^{\mathrm{init}}_{C,a} = w_a^\top u_{g(a)} + b_a,
\label{eq:branchC_logit}
\end{equation}
and a co-resistance refinement lets one drug's prediction modulate a related one's,
\begin{equation}
r = \mathrm{MLP}\!\left(\left[z_C \, ; \, \ell^{\mathrm{init}}_{C}\right]\right),
\qquad
\ell_C = \ell^{\mathrm{init}}_{C} + \alpha\, r,
\label{eq:branchC_cores}
\end{equation}
with $\alpha$ a learned scalar initialised at 0.1, so the model ignores the step until the data
support it and no spurious drug-to-drug relation is baked in at the start. The shared trunk is what
multi-task learning is for; the per-class branches limit how far a misleading pattern in the trunk
can propagate.

\begin{figure}[t]\centering
\figbox{Figure S1. Co-resistance}{6cm}{Pairwise co-resistance ($\phi$) between the seven
antibiotics, each cell computed on the strains labelled for both drugs, with $n$ shown in the cell.
Lines separate drug classes. The two cephalosporins reach 0.84; cross-class pairs reach
0.55.}{results/data\_inventory/strain\_table.csv}
\caption{\textbf{Co-resistance is only partly explained by drug class.} \todo{caption}}
\label{fig:coresistance}
\end{figure}

\subsection{The mixture-of-experts gate}
Averaging the three branches would assume they are equally useful, and they are not: where
resistance comes from an acquired gene, branch B should dominate and branch A has almost nothing to
read, and the reverse holds for a point mutation in a target enzyme. A fixed average drowns the
stronger opinion in two weaker ones. The gate \cite{jacobs1991moe,shazeer2017moe} therefore learns
the weights, per strain and per drug. The three summaries are concatenated into
$z = [z_A; z_B; z_C] \in \mathbb{R}^{192}$ and passed through a small network and a softmax,
\begin{equation}
g = \mathrm{softmax}\big(\mathrm{MLP}([z_A; z_B; z_C])\big), \qquad g \in \mathbb{R}^{B \times 7 \times 3},
\label{eq:moe_gate}
\end{equation}
so that for each drug the three weights are non-negative and sum to one, and the fused score is
\begin{equation}
\ell^{\mathrm{fin}}_{a} = \sum_{e=1}^{3} g_{a,e}\,\ell_{e,a},
\qquad
p^{\mathrm{neu}}_{a} = \sigma\!\left(\ell^{\mathrm{fin}}_{a}\right).
\label{eq:fusion}
\end{equation}

Such a gate has a characteristic failure mode. Early in training one branch is slightly ahead by
chance; the gate shifts weight to it; the larger gradient makes it learn faster; after a few epochs
the other two have stopped training and the model has silently become a single-branch network. We
measure the gate's spread for each drug by its Shannon entropy,
\begin{equation}
H(g_a) = -\sum_{e=1}^{3} g_{a,e} \log g_{a,e},
\label{eq:gate_entropy}
\end{equation}
largest when the three weights are equal and smallest when the gate has collapsed, and add to the
loss
\begin{equation}
\mathcal{L}_{\mathrm{ent}}
  = -\frac{\lambda_{\mathrm{ent}}}{7} \sum_{a=1}^{7} H(g_a),
\label{eq:entropy_reg}
\end{equation}
with $\lambda_{\mathrm{ent}} = 0.01$, which is minimised by a gate that spreads its weight while
leaving it free to favour one branch where the evidence supports it. We track the mean weight of
each branch during training and raise $\lambda_{\mathrm{ent}}$ automatically if any branch falls
below a threshold. Each branch also carries an auxiliary loss on its own scores, without which a
branch the gate currently ignores receives almost no gradient and stops improving. The objective is
therefore binary cross-entropy on the fused scores, plus \eqref{eq:entropy_reg}, plus the auxiliary
terms, computed over observed labels only.

\subsection{The stacking meta-learner}
The decision is made by a per-drug logistic regression over the neural probability and the five
classical out-of-fold probabilities, which is stacked generalisation in Wolpert's sense
\cite{wolpert1992stacked}. For drug $a$,
\begin{equation}
\phi_a = \left[\,p^{\mathrm{neu}}_a,\; p^{(1)}_a,\; \dots,\; p^{(5)}_a\,\right],
\qquad
\eta_a = w_a^\top \phi_a + b_a,
\qquad
\hat{p}_a = \sigma(\eta_a),
\label{eq:meta}
\end{equation}
with an $L_2$ penalty and $C = 1$, and the call is $\hat{y}_a = \mathbf{1}[\hat{p}_a > 0.5]$. One
meta-learner per drug is deliberate: the drugs rely on different base learners, and a shared
meta-learner would force one trade-off on all seven. Unlike a fixed average, the weights are
learned, which is the difference between averaging six opinions and weighting six opinions by their
track records.

A neural meta-learner was also tested: a small feed-forward layer producing one token per drug, a
four-head self-attention layer letting the seven tokens exchange information, and a separate
decoding head per drug. It was aimed at exactly the co-resistance structure a per-drug linear model
cannot represent, and it scored lower --- 0.825 against 0.850 average precision on the random split.
The input dimension is not the problem; folding the neural embedding into the meta-features widens
the vector to 170 dimensions, but only about 5,000 strains sit behind it, and an attention layer
with seven decoding heads fits noise at that size. The linear meta-learner is also directly
readable, which the deeper one is not.

\section{Supplementary Methods: the presence/absence matrix repair}

The release encodes the 32,441 pan-genome loci as 0 and 2 without documenting which value means
present, and the two readings are biologically opposite. Resolving it by an external check rather
than a guess: under one reading the median locus prevalence is 0.966 and the average strain carries
about 28,700 loci; under the other the median is 0.034 and the average strain carries 3,709 (median
3,349, range 2,141--8,333). An \Ecoli{} genome holds roughly 4,000--5,000 genes, which selects the
second. Note that the two statistics are not derivable from one another --- the median is over loci
and the count over strains, and the prevalence distribution is strongly skewed, so the count is not
the median times 32,441.

Rebuilding in that orientation flipped 10,511 loci and left 21,930 unchanged, with no locus
inconsistent; 2,924 loci are monomorphic, and no cell is missing or heterozygous. Three checks
support the result: the core genes \textit{gyrA}, \textit{parC}, \textit{rpoB} and \textit{recA}
reach prevalence 1.0, where the original orientation capped every locus at 0.504; locus prevalence
correlates with ECOR reference frequencies at $\rho = 0.689$; and of 3,643 loci expected to be core,
0.1\% fall below 50\% prevalence. The per-strain locus count is now checked automatically on every
matrix build, because the wrong reading raises no error --- the model would train normally with
every gene inverted, and since this view carries most of the gate's weight, everything downstream
would have been wrong in a way nothing would have revealed.

Of the 6,962 strains, 6,949 appear in the PLINK sample list. The remaining 13 (0.19\%) are carried
as all-zero vectors, which conflates a missing record with an empty accessory genome; the correct
handling is to extend the label mask to this view for those rows. At 0.19\% of the collection,
distributed across partitions in proportion, the effect is below the resolution of the intervals
reported, so we record the defect rather than publishing a corrected set of results.

A separate repair concerns the SNP matrix. A contiguous block of 95 columns (indices
169,213--169,307, spanning two genes) held 90,436 cells outside the valid range 0--3; 117,598 cells
across 48 columns were rewritten, all inside that block, after which the matrix range is 0--3.

\section{Supplementary Methods: variant filtering and gene selection}

Table~\ref{tab:variants} records the variant-level filtering. We read 2,451,184 variants, matching
the figure the authoring group reports, which confirmed the file was being read correctly before
anything was built on it. The frequency filter alone removed 76\%, most of them present in two or
three strains; a pattern seen in so few strains cannot be expected to hold in a new one.

\begin{table}[t]\centering\small
\caption{Variant filtering during construction of the SNP matrix.}
\label{tab:variants}
\begin{tabular}{lr}
\toprule
Stage & Variants \\
\midrule
Streamed from the variant file & 2,451,184 \\
Removed: insertion or deletion & 43,799 \\
Removed: minor allele frequency below 1\% & 1,871,946 \\
\midrule
Kept & 535,439 \\
\bottomrule
\end{tabular}
\end{table}

Gene-level reduction was a modelling necessity rather than an optimisation: 535,439 features over
6,962 strains is about 77 candidate features per strain. Because the model reads the genome at the
gene level, reduction is at the gene level too. The variability filter keeps genes with at least
five SNPs in the training partition of the split being fitted (18,975 $\rightarrow$ 13,799 on the
random split), since fewer variable positions cannot support a within-gene pattern. The importance
filter then ranks genes by summed gradient-boosting gain over the seven per-drug models, each
trained on that drug's top 50,000 $\chi^2$-filtered SNPs, and keeps the top 5,000. Both filters are
refitted from scratch for each split; had test labels been visible during ranking, they would have
changed which features exist and every score would have been optimistic by an unknown amount.

The retained 5,000 genes carry 262,225 SNPs, 49.0\% of the filtered variants and 10.7\% of those
originally streamed. Gene sizes vary widely --- smallest 5 SNPs, median 32, mean 52, largest 717 ---
which fixed the block length at 64: about 28\% of selected genes are truncated and the rest padded
and masked, with an average of 35 real positions per block. A longer block would waste memory on
most genes and a shorter one would cut the median gene.

\section{Supplementary Results}

\subsection{Threshold tuning on the internal splits}
Table~\ref{tab:thresholds} gives the per-drug thresholds fitted on the validation partition and
their effect on the random-split test partition. Neither objective rescues the clinical numbers:
maximising $F_2$ lowers the mean very major error rate from 0.241 to 0.133 while raising the mean
major error rate from 0.055 to 0.141, and constraining the major error rate to the 3\% benchmark
pushes the very major error rate to 0.308. No drug meets both limits at any threshold tested.

\begin{table}[t]\centering\small
\caption{Per-drug thresholds fitted on validation and applied to the random-split test partition.
VME, very major error; ME, major error.}
\label{tab:thresholds}
\begin{tabular}{lrrrrrrrr}
\toprule
 & \multicolumn{2}{c}{At 0.5} & \multicolumn{3}{c}{Maximising $F_2$} & \multicolumn{3}{c}{ME capped at 3\%} \\
Drug & VME & ME & thr. & VME & ME & thr. & VME & ME \\
\midrule
AMP & 0.107 & 0.099 & 0.16 & 0.044 & 0.255 & 0.88 & 0.363 & 0.022 \\
CIP & 0.071 & 0.015 & 0.13 & 0.045 & 0.029 & 0.10 & 0.045 & 0.035 \\
AMC & 0.394 & 0.105 & 0.12 & 0.080 & 0.418 & 0.77 & 0.669 & 0.033 \\
SXT & 0.218 & 0.086 & 0.37 & 0.182 & 0.114 & 0.77 & 0.412 & 0.028 \\
GEN & 0.415 & 0.023 & 0.22 & 0.319 & 0.049 & 0.21 & 0.319 & 0.051 \\
CTX & 0.190 & 0.026 & 0.11 & 0.095 & 0.050 & 0.18 & 0.124 & 0.041 \\
CAZ & 0.291 & 0.034 & 0.10 & 0.164 & 0.075 & 0.30 & 0.227 & 0.047 \\
\midrule
Mean & 0.241 & 0.055 & --- & 0.133 & 0.141 & --- & 0.308 & 0.037 \\
\bottomrule
\end{tabular}
\end{table}

\subsection{Near-duplicate removal on the random split}
Removing the 47 test strains within 0.5\% SNP divergence of a training strain changed the mean
average precision from 0.850 to 0.846 ($-0.005$). Per drug: AMP $-0.000$, CIP $-0.002$,
AMC $+0.008$, SXT $+0.000$, GEN $-0.003$, CTX $-0.012$, CAZ $-0.023$. The effect is small and not
directional. On the lineage holdout no test strain was flagged, so the deduplicated and raw holdout
numbers are identical.

\subsection{Split integrity checks}
All three splits are disjoint and complete over the 6,962 strains. On the lineage holdout, 240
training and 103 test sequence types share none; validation and test share none; and all 475 strains
without a sequence type are outside the test partition. For each split, the saved predictions of
every run reported in the main text were verified to carry that split's exact training, validation
and test indices in the same order, and the neural and classical arms were verified to cover the
same strains in the same order before stacking. No test genome is identical to a training genome on
the filtered SNPs in any split.

\subsection{Phylogroup composition}
B2 accounts for 3,542 strains (50.9\%), unassigned 1,697 (24.4\%), D 723 (10.4\%), A 317 (4.6\%),
F 307 (4.4\%), B1 258 (3.7\%), C 114 (1.6\%) and E 4 (0.06\%). Three properties limit what can be
read from any phylogroup analysis: B2 alone is over half the collection, so a phylogroup split is
largely B2 against the rest; nearly a quarter of strains have no assignment, so the unassigned group
is not a group; and phylogroup E has four members, too few for any per-group estimate. On the
phylogroup-stratified split, a 70/15/15 allocation places at most one E strain in the test
partition, so that group is present by construction and carries no evidential weight.

\subsection{Per-drug results for the pre-registered experiments}
\todo{Deposit the per-drug tables for W6--W14 here, or cite the repository paths; the
pre-registration documents with their dated amendments go in this supplement as well.}

\bibliographystyle{unsrtnat}
\bibliography{refs}

\end{document}
